\lecture{5}{Fri 09 Sep 2022 14:33}{Relativistic Dynamics}

\subsection{Relativistic Dynamics}

\begin{example}(Classical Collision)
	In some reference frame $O'$:
	 \begin{align*}
		m_1 &=  2m \\
		v_1 &=  0 \\
		m_2 &= m \\
		v_2 &= -0.75 c \hat{x} 
	.\end{align*}

\begin{figure}[ht]
    \centering
    \incfig{classical-collision}
    \caption{classical-collision}
    \label{fig:classical-collision}
\end{figure}

Apply conservation of momentum:
\begin{align*}
	p_i' &= m_1v_{1i}'+m_2 v_{2i}' = -0.75 m c \\
	p_f' &=  -m_1 v_{1f}' + m_2 v_{2f}'  = -0.75mc
.\end{align*}
Assume elastic collision and apply conservaton of energy:
\[
	\frac{1}{2} m_1 v_{1i}'^2 + 
\frac{1}{2} m_2 v_{2i}'^2 = 
\frac{1}{2} m_1 v_{1f}'^2 + 
\frac{1}{2} m_2 v_{2f}'^2 
.\] 

Solutions:
\begin{align*}
	v_{1f}'&=-0.5c\\
	v_{2f}'&=0.25c
.\end{align*}

\end{example}

\begin{problem}
	If we apply special relativity to transform the velocities into a different reference frame, we do not conserve momentum and energy.
\end{problem}

\begin{definition}[Relativistic momentum]
	\[
		\vec{p} = \gamma m \vec{v} 
	.\] 
\end{definition}
Using this definition, we conserve momentum in both frames, but do not get the same momentum in both frames.


\begin{theorem}[Work-Energy Theorem]
	\[
		W = \int F \cdot d x = \int \frac{\partial p}{\partial t} dx = \int \frac{\partial x}{\partial t} dp = \int v dp = pv - \int p dv
	.\] 
\end{theorem}

Assume this works relativistically:

 \begin{align*}
	 K &= (\gamma mv)v - \int_0^v (\gamma mv) dv\\
	 &= \gamma m v^2 + \gamma^{ -1 } m c^2 - m c^2 \\
	 &= \gamma m c^2 - m c^2
.\end{align*} 

\begin{definition}[Relativistic Kinetic Energy]
	\[
		K = \gamma m c^2 - m c^2
	\]
	where we define
	\[
		E = \gamma m c^2
	.\]  
	to be the \textit{relativistic total energy}, and define \[
		E_0 = m c^2
	\] to be the \textit{relativistic rest energy}.

\end{definition}
The \textit{total energy} is conserved under frame transformation. 
We have seen that kinetic energy is a relativistic effect.


Let $v\ll c$ and do Taylor Expansion:
\[
	K \approx m c^2 (1 + \frac{1}{2}\frac{v^2}{c^2}) - m c^2 = \frac{1}{2} m v^2
.\]
So this reduces into Newton Mechanics at low speed.

\begin{remark}
	Alternative forms: 
	\begin{itemize}
		\item $E^2 = (pc)^2 + \left(mc^2\right)^2$
		\item $E^2 = p^2 + m^2$ when we take $c=1$ 
		\item $E = pc$ when we take $m=0$
	\end{itemize}
\end{remark}

\begin{remark}(Units)
	proton has rest energy of  \SI{938}{\mega\electronvolt}. It has mass \SI{938}{\mega\electronvolt\per c \squared}. The unit for momentum is \SI{}{\mega \electronvolt \per c}. \color{blue} Pick the unit that make everything look nice!
\end{remark}


